\paragraph{Signed conditional slices at genuine theta zeros.}
\label{sec:thetasignedslices}

The division-free Duhamel formula identifies positivity
of the complete weighted Fourier correlation as the
unproved step toward RH. One tempting strengthening
is to insist that the Fourier transform of each
fixed-difference slice be nonnegative. The following
argument tests this proposal \emph{within the genuine
Jacobi-theta family}, without assuming RH.

For any heat time $t$ write
\[
\nu_t(u)=\tfrac12e^{tu^2}\Phi(u),\qquad
H_t(x)=\int_{\mathbb R}\nu_t(u)e^{ixu}\,du .
\]
For $b,x\in\mathbb R$ define the real, even
fixed-difference slice
\begin{equation}
C_t(b,x)=\frac12\int_{\mathbb R}
\nu_t\!\left(\frac{w+b}{2}\right)
\nu_t\!\left(\frac{w-b}{2}\right)
\cos(xw)\,dw .
\label{eq:thetafixedslice}
\end{equation}
The two variables have a distinct meaning:
$w$ is the sum of the original pair of
Fourier variables, and $b$ their difference.
This is an ambiguity-function-type slice;
it is not the positive density obtained by
integrating \emph{without} the cosine factor.

\begin{proposition}[Exact slice transform and moments]
\label{prop:thetafixedsliceidentities}
For every real $t,x,\eta$ and integer $n\ge0$,
\begin{align}
\int_{\mathbb R}C_t(b,x)e^{i\eta b}\,db
&=H_t(x+\eta)H_t(x-\eta),\label{eq:slicesymbol}\\
\frac{1}{(2n)!}
\int_{\mathbb R}b^{2n}C_t(b,x)\,db
&=L_{n,t}(x).\label{eq:slicemoments}
\end{align}
Both identities are absolutely convergent.
\end{proposition}

\begin{proof}
In the double integral make the linear substitution
$u=(w+b)/2$, $v=(w-b)/2$, for which
$du\,dv=\tfrac12\,dw\,db$. The transform
with $\cos(xw)e^{i\eta b}$ then factors into
the average of
$H_t(x+\eta)H_t(x-\eta)$ and
$H_t(-x+\eta)H_t(-x-\eta)$; these coincide
because $\nu_t$ is even and real.
This proves \eqref{eq:slicesymbol}.
Differentiation of the latter identity at
$\eta=0$ gives the even moments.
Equivalently,
\[
H_t(x+iy)H_t(x-iy)
=\sum_{n\ge0}L_{n,t}(x)y^{2n}
\]
and the change $y=i\eta$ matches
the characteristic function of $C_t$.
\end{proof}

\begin{theorem}[Necessary sign compensation at every
simple real theta zero]\label{thm:thetaslicesign}
Suppose $H_t(a)=0$ and $H_t'(a)\ne0$
for a real $a$. Then the genuine theta
slice $b\mapsto C_t(b,a)$ takes
\emph{both positive and negative values}.
More precisely,
\begin{equation}
\int C_t(b,a)\,db=0,\qquad
\int b^2C_t(b,a)\,db=2H_t'(a)^2>0 .
\label{eq:twosignmoment}
\end{equation}
Its Fourier transform in $b$ satisfies
\begin{equation}
\widehat C_t(\eta,a)
= -H_t'(a)^2\eta^2+O(\eta^4)<0
\label{eq:negativeslicesymbol}
\end{equation}
for all sufficiently small nonzero real $\eta$.
Consequently neither a pointwise nonnegative
slice decomposition nor a positive
convolution-square (positive-definite)
factorization of these \emph{individual}
fixed-difference slices can hold.

These conclusions apply unconditionally
to the actual theta heat family at $t=0$
using any of its established simple real
zeros, and at $t=1/2$ by the classical
real-simple-zero theorem.
\end{theorem}
\begin{proof}
At a simple real zero $a$,
$L_{0,t}(a)=H_t(a)^2=0$, while
$L_{1,t}(a)=H_t'(a)^2>0$.
Proposition~\ref{prop:thetafixedsliceidentities}
yields \eqref{eq:twosignmoment}.
If $C_t(\cdot,a)$ were nonnegative
almost everywhere, its zero integral
would force it to vanish almost
everywhere, contradicting its strictly
positive second moment. Since $C_t$ is
continuous, real and integrable, it must
take both signs.

Taylor's theorem and reality on the real
axis give
\[
H_t(a+\eta)H_t(a-\eta)
=-H_t'(a)^2\eta^2+O(\eta^4).
\]
Equation~\eqref{eq:slicesymbol} proves
\eqref{eq:negativeslicesymbol}.
A positive-definite integrable function
has a nonnegative Fourier transform;
equivalently any sum or positive mixture
of convolution squares has nonnegative
Fourier symbol. The strict negative
value contradicts such a representation.
The stated actual-theta applications use
known real simple zeros, not an assumption
that \emph{all} theta zeros are real.
\end{proof}

\paragraph{Two separate tests.}
An exact rational control uses the
even two-atom measure
$\nu=\tfrac12(\delta_{-1}+\delta_1)$,
for which $H(x)=\cos x$.
At its simple zero $a=\pi/2$,
the signed slice has masses
\[
C(-2,a)=\tfrac14,\qquad
C(0,a)=-\tfrac12,\qquad
C(2,a)=\tfrac14
\]
in the atomic sense. Their total
is zero, whereas their second moment
is $2$. This normalization test is
checked by exact rational arithmetic.

For the \emph{actual} theta series,
a separate double-precision Simpson
diagnostic at $x=28.269$ gives
\[
C_0(0,x)\simeq-3.26794\cdot10^{-7},
\quad
C_0(\tfrac12,x)\simeq+9.74345\cdot10^{-8}.
\]
The signs were reproduced using
8,000 and 16,000 panels. These sampled
signs are not interval enclosures and
are not used in the theorem.
The code is
\path{certificates/xi_theta_conditional_slice_attack.py}.
The algebraic signs are also
formally checked in
\path{XiKernel/ThetaConditionalSliceObstruction.lean};
the infinite theta Fourier integral
is not formalized there.

\paragraph{The remaining viable decomposition.}
This result invalidates positivity
\emph{before} integration in the
difference variable $b$; it does not
refute positivity of the completed
Laguerre correlations
\[
L_{n,t}(x)=\frac1{(2n)!}
\int b^{2n}C_t(b,x)\,db.
\]
The weights $b^{2n}$ can compensate
negative and positive slices,
as \eqref{eq:twosignmoment} already
demonstrates at a simple zero.
To obtain a non-circular RH proof one
would need a theta-specific inequality
for the \emph{fully weighted} signed
integral simultaneously for all
$n,t,x$, or a different complex
zero-exclusion theorem. Showing
positivity of every fixed slice,
or positive definiteness in the
difference variable, is not that
missing mechanism.

The correlation and Fourier-positivity
setting has significant prior art
in the work of Csordas~\cite{Cs}
and Dimitrov--Xu~\cite{DX}.
The identities above are deliberately
presented as an obstruction to the
proposed proof mechanism, not a
new criterion for RH or a proof of it.
